\documentclass[a4paper,12pt]{article}

\usepackage[T2A]{fontenc}
\usepackage[utf8]{inputenc}
\usepackage[russian]{babel}
\usepackage{amsmath,amssymb,amsthm,mathtools}
\usepackage{helvet}
\renewcommand{\familydefault}{\sfdefault}
\usepackage{icomma}
\usepackage[a4paper,margin=20mm,headheight=15pt]{geometry}
\usepackage{microtype}
\usepackage{tikz}
\usetikzlibrary{calc,arrows.meta,positioning,shapes.geometric}
\usepackage{tkz-euclide}
\usepackage{pgfplots}
\pgfplotsset{compat=1.18}
\usepackage{tabularx,booktabs,longtable}
\usepackage{enumitem}
\usepackage{hyperref}
\usepackage{parskip}
\usepackage{fancyhdr}
\usepackage{setspace}
\usepackage{xcolor}

\definecolor{accent}{HTML}{008080}
\definecolor{darkaccent}{HTML}{004D4D}
\definecolor{textgray}{HTML}{333333}
\definecolor{softgray}{HTML}{F2F4F4}

\hypersetup{
  colorlinks=true,
  linkcolor=darkaccent,
  urlcolor=darkaccent
}

\onehalfspacing
\setlength{\parindent}{0pt}

\setlist[itemize]{
  leftmargin=1.5em,
  itemsep=0.25em,
  topsep=0.3em
}

\setlist[enumerate]{
  leftmargin=1.7em,
  itemsep=0.45em,
  topsep=0.35em
}

\pagestyle{fancy}
\fancyhf{}
\fancyhead[L]{\small\color{textgray}Трапеция: ключевые методы}
\fancyhead[R]{\small\color{textgray}София Кальней}
\fancyfoot[C]{\thepage}
\renewcommand{\headrulewidth}{0.3pt}

\newcommand{\sectionline}{
  \par
  \vspace{0.2em}
  {\color{accent}\hrule height 0.8pt}
  \vspace{0.7em}
}

\newcommand{\important}[1]{
  \textbf{\color{darkaccent}#1}
}

\newcommand{\formula}[1]{
  \[
  \boxed{#1}
  \]
}

\tikzset{
  flowstart/.style={
    ellipse,
    draw=darkaccent,
    line width=0.8pt,
    align=center,
    minimum width=3.8cm,
    minimum height=0.95cm,
    fill=softgray,
    font=\small
  },
  flowaction/.style={
    rounded corners=3pt,
    draw=accent,
    line width=0.8pt,
    align=center,
    text width=5.9cm,
    minimum height=0.9cm,
    font=\small
  },
  flowdecision/.style={
    diamond,
    aspect=2.5,
    draw=darkaccent,
    line width=0.8pt,
    align=center,
    text width=4.4cm,
    inner sep=0.5pt,
    font=\small
  },
  flowarrow/.style={
    -{Latex[length=3mm]},
    line width=0.8pt,
    draw=darkaccent
  }
}

\begin{document}

% =========================================================
% Титульный лист
% =========================================================

\begin{titlepage}
\thispagestyle{empty}
\centering

\vspace*{2.6cm}

{\Large\color{darkaccent}\textbf{Чек-лист}}\\[0.8em]

{\LARGE\textbf{
Трапеция: ключевые методы, свойства\\
и дополнительные построения
}}\\[1.2em]

{\large ОГЭ по математике}\\[2.2em]

{\large \today}\\[2.5em]

{\large Лёвин Артём Александрович}\\[0.35em]

{\large\textbf{эксклюзивно для Софии Кальней}}\\[3em]

\begin{minipage}{0.78\textwidth}
\centering
\color{textgray}
Главная идея занятия: в задачах о трапеции целенаправленно искать
треугольники, к которым можно применить уже известные теоремы
и формулы.
\end{minipage}

\vfill

\begin{tikzpicture}[remember picture,overlay]
\draw[
  accent!65,
  line width=1.2pt
]
  ($(current page.south west)+(18mm,18mm)$)
  .. controls
  ($(current page.south west)+(65mm,33mm)$)
  and
  ($(current page.south east)+(-70mm,7mm)$)
  ..
  ($(current page.south east)+(-18mm,19mm)$);
\end{tikzpicture}

\end{titlepage}

% =========================================================
% Основной чек-лист
% =========================================================

\section*{1. Базовые формулы трапеции}
\sectionline

Пусть основания трапеции равны $a$ и $b$, причём $a<b$,
а высота равна $h$.

Площадь трапеции:

\formula{
S=\frac{a+b}{2}\,h
}

Средняя линия трапеции:

\formula{
m=\frac{a+b}{2}
}

Поэтому

\[
S=mh.
\]

\important{Ключевая ситуация.}
Основания известны, а высота отсутствует. Тогда основной вопрос задачи:
как построить треугольник, в котором эту высоту удастся вычислить?

% ---------------------------------------------------------

\section*{2. Главный приём: параллельный перенос}
\sectionline

Параллельный перенос особенно полезен тогда, когда после дополнительного
построения возникает треугольник с достаточным количеством известных
элементов.

\subsection*{2.1. Перенос боковой стороны}

Пусть в трапеции $ABCD$

\[
AD\parallel BC,
\qquad
AD=b,
\qquad
BC=a,
\]

а боковые стороны равны

\[
AB=c,
\qquad
CD=d.
\]

Через вершину $C$ проведём $CF\parallel AB$ до пересечения с основанием
$AD$.

\begin{center}
\begin{tikzpicture}[scale=0.8]

\tkzSetUpLine[
  line width=1pt,
  color=accent
]

\tkzSetUpPoint[
  color=accent,
  fill=accent,
  size=3
]

\tkzSetUpLabel[
  color=black,
  font=\small
]

\tkzDefPoint(0,0){A}
\tkzDefPoint(7,0){D}
\tkzDefPoint(1.2,3){B}
\tkzDefPoint(4.2,3){C}
\tkzDefPoint(3,0){F}

\tkzDrawPolygon(A,B,C,D)
\tkzDrawSegment(C,F)

\tkzDrawPoints(A,B,C,D,F)

\tkzLabelPoints[below](A,F,D)
\tkzLabelPoints[above](B,C)

\node[below] at (1.5,0) {$AF=a$};
\node[below] at (5.0,0) {$FD=b-a$};
\node[above] at (2.7,3) {$BC=a$};

\end{tikzpicture}
\end{center}

Четырёхугольник $ABCF$ является параллелограммом. Следовательно,

\[
CF=AB=c,
\qquad
AF=BC=a.
\]

Поэтому

\[
FD=AD-AF=b-a.
\]

Получаем треугольник $CFD$ со сторонами

\[
c,\qquad d,\qquad b-a.
\]

Высота этого треугольника к стороне $FD$ совпадает с высотой исходной
трапеции.

Если все три стороны полученного треугольника известны, удобно применить
формулу Герона:

\[
p=\frac{x+y+z}{2},
\]

\[
S_{\triangle}
=
\sqrt{
p(p-x)(p-y)(p-z)
}.
\]

С другой стороны,

\[
S_{\triangle}
=
\frac12(b-a)h.
\]

Приравнивая эти выражения, находим высоту $h$.

\subsubsection*{Пример}

Основания трапеции равны $16$ и $44$, боковые стороны --- $17$ и $25$.

После параллельного переноса получаем треугольник со сторонами

\[
17,\qquad 25,\qquad 44-16=28.
\]

Его полупериметр:

\[
p=\frac{17+25+28}{2}=35.
\]

По формуле Герона:

\[
S_{\triangle}
=
\sqrt{
35\cdot18\cdot10\cdot7
}
=
210.
\]

С другой стороны,

\[
S_{\triangle}
=
\frac12\cdot28\cdot h.
\]

Следовательно,

\[
14h=210,
\qquad
h=15.
\]

Площадь исходной трапеции:

\[
S
=
\frac{16+44}{2}\cdot15
=
450.
\]

\subsection*{2.2. Перенос диагонали}

Параллельно переносить можно также диагональ.

Если основания трапеции равны $a$ и $b$, а диагонали имеют длины
$p$ и $q$, после подходящего параллельного переноса возникает треугольник
со сторонами

\[
p,\qquad q,\qquad a+b.
\]

Высота такого треугольника к стороне $a+b$ совпадает с высотой трапеции.

\important{Критерий удачного построения.}
После дополнительного построения должен появиться треугольник, для которого
известно достаточно данных, чтобы вычислить его высоту или площадь.

% =========================================================
% Блок-схема
% =========================================================

\clearpage
\thispagestyle{empty}

\begin{center}
{\Large\color{darkaccent}\textbf{
Алгоритм: как искать высоту трапеции
}}
\end{center}

\vspace{1.4em}

\begin{center}
\begin{tikzpicture}[node distance=0.55cm]

\node[flowstart] (start) {
Даны элементы трапеции;\\
нужно найти площадь
};

\node[flowaction,below=of start] (formula) {
Запишите
$
S=\dfrac{a+b}{2}h
$.\\
Определите, что требуется найти $h$.
};

\node[
  flowdecision,
  below=0.75cm of formula
] (triangle) {
Есть удобный треугольник,\\
содержащий высоту?
};

\node[
  flowaction,
  below=0.75cm of triangle
] (build) {
Если нет, проведите высоту или выполните\\
параллельный перенос боковой стороны/диагонали.
};

\node[
  flowdecision,
  below=0.75cm of build
] (known) {
В полученном треугольнике\\
достаточно известных данных?
};

\node[
  flowaction,
  below=0.75cm of known
] (calc) {
Найдите его площадь или высоту:\\
Герон, Пифагор, свойства $30^\circ$,\\
равнобедренного или прямоугольного треугольника.
};

\node[
  flowaction,
  below=of calc
] (back) {
Подставьте найденную высоту\\
в формулу площади трапеции.
};

\node[
  flowstart,
  below=of back
] (finish) {
Вычислите $S$ и проверьте результат.
};

\draw[flowarrow]
(start)--(formula);

\draw[flowarrow]
(formula)--(triangle);

\draw[flowarrow]
(triangle)
--
node[right]{\small нет / требуется построение}
(build);

\draw[flowarrow]
(build)--(known);

\draw[flowarrow]
(known)
--
node[right]{\small да}
(calc);

\draw[flowarrow]
(triangle.east)
--
++(2.2,0)
|-
node[pos=0.25,right]{\small да}
(calc.east);

\draw[flowarrow]
(known.west)
--
++(-2.2,0)
|-
node[pos=0.25,left]{\small нет}
(build.west);

\draw[flowarrow]
(calc)--(back);

\draw[flowarrow]
(back)--(finish);

\end{tikzpicture}
\end{center}

\clearpage

% =========================================================

\section*{3. Метод равенства площадей}
\sectionline

Если высота входит в одну формулу площади, а саму площадь можно найти
другим способом, полезно вычислить одну и ту же площадь двумя способами
и приравнять результаты.

Пусть треугольник имеет стороны $x$, $y$, $z$, а $h$ --- высота
к стороне $z$.

Тогда

\[
S_{\triangle}
=
\frac12zh.
\]

По формуле Герона:

\[
S_{\triangle}
=
\sqrt{
p(p-x)(p-y)(p-z)
},
\]

где

\[
p=\frac{x+y+z}{2}.
\]

Следовательно,

\[
\frac12zh
=
\sqrt{
p(p-x)(p-y)(p-z)
}.
\]

Отсюда

\[
h
=
\frac{
2\sqrt{
p(p-x)(p-y)(p-z)
}
}{z}.
\]

\important{Контроль.}
Формулу Герона применяем только после того, как известны все три стороны
рассматриваемого треугольника.

% =========================================================

\section*{4. Равнобедренная трапеция: проекции и средняя линия}
\sectionline

Пусть $ABCD$ --- равнобедренная трапеция,

\[
AD=b,
\qquad
BC=a,
\qquad
b>a,
\]

и

\[
CH\perp AD.
\]

\begin{center}
\begin{tikzpicture}[scale=0.85]

\tkzSetUpLine[
  line width=1pt,
  color=accent
]

\tkzSetUpPoint[
  color=accent,
  fill=accent,
  size=3
]

\tkzSetUpLabel[
  color=black,
  font=\small
]

\tkzDefPoint(0,0){A}
\tkzDefPoint(8,0){D}
\tkzDefPoint(2,3){B}
\tkzDefPoint(6,3){C}
\tkzDefPoint(6,0){H}

\tkzDrawPolygon(A,B,C,D)
\tkzDrawSegment(C,H)

\tkzMarkRightAngle[size=0.28](C,H,D)

\tkzDrawPoints(A,B,C,D,H)

\tkzLabelPoints[below](A,H,D)
\tkzLabelPoints[above](B,C)

\node[above] at (4,3) {$a$};
\node[below] at (4,0) {$b$};

\end{tikzpicture}
\end{center}

Для равнобедренной трапеции получаем

\[
DH=\frac{b-a}{2},
\]

а большая часть нижнего основания равна

\[
AH
=
b-DH
=
b-\frac{b-a}{2}
=
\frac{a+b}{2}.
\]

Следовательно,

\formula{
AH=m=\frac{a+b}{2}
}

Итак:

\[
\boxed{
\text{большой отрезок}
=
\frac{a+b}{2}
}
\]

и

\[
\boxed{
\text{малый отрезок}
=
\frac{b-a}{2}
}.
\]

Эти две формулы особенно полезны при проведении высоты в равнобедренной
трапеции.

% ---------------------------------------------------------

\subsection*{4.1. Диагональ под углом $60^\circ$}

Пусть диагональ равнобедренной трапеции равна $d$ и образует с основанием
угол $60^\circ$.

Её проекция на основание равна

\[
d\cos60^\circ
=
\frac d2.
\]

Для равнобедренной трапеции эта проекция равна

\[
\frac{a+b}{2}.
\]

Поэтому средняя линия равна

\formula{
m=\frac d2
}

Например, при $d=10$

\[
m=5.
\]

% ---------------------------------------------------------

\subsection*{4.2. Перпендикулярные диагонали}

У равнобедренной трапеции диагонали равны.

Если дополнительно известно, что диагонали перпендикулярны, после
параллельного переноса одной диагонали получается прямоугольный
равнобедренный треугольник.

Его гипотенуза равна сумме оснований:

\[
a+b.
\]

Высота исходной трапеции становится медианой к гипотенузе полученного
прямоугольного треугольника.

В прямоугольном треугольнике медиана, проведённая к гипотенузе,
равна половине гипотенузы. Следовательно,

\[
h
=
\frac{a+b}{2}.
\]

Но

\[
m
=
\frac{a+b}{2}.
\]

Поэтому

\formula{
h=m
}

Для площади получаем особенно короткую формулу:

\formula{
S=m^2
}

\important{Важно.}
Перпендикулярность диагоналей не является общим свойством всех
равнобедренных трапеций. Она должна следовать из условия конкретной задачи.

% =========================================================

\section*{5. Вписанная окружность в трапецию}
\sectionline

Пусть окружность с центром $O$ вписана в трапецию $ABCD$, где

\[
AD\parallel BC.
\]

Центр вписанной окружности лежит на биссектрисах углов.

Следовательно,

\[
\angle BAO
=
\frac12\angle BAD,
\]

\[
\angle ABO
=
\frac12\angle ABC.
\]

Поскольку $AD\parallel BC$, углы

\[
\angle BAD
\quad\text{и}\quad
\angle ABC
\]

являются односторонними при боковой стороне $AB$. Поэтому

\[
\angle BAD+\angle ABC=180^\circ.
\]

Рассмотрим треугольник $AOB$:

\[
\angle AOB
=
180^\circ
-
\angle BAO
-
\angle ABO.
\]

Подставляем выражения через половины углов:

\[
\angle AOB
=
180^\circ
-
\frac12
\left(
\angle BAD+\angle ABC
\right).
\]

Следовательно,

\[
\angle AOB
=
180^\circ-90^\circ
=
90^\circ.
\]

\formula{
\angle AOB=90^\circ
}

Таким образом,

\[
\triangle AOB
\]

является прямоугольным.

% ---------------------------------------------------------

\subsection*{5.1. Точка касания на боковой стороне}

Пусть окружность касается стороны $AB$ в точке $T$.

Радиус, проведённый в точку касания, перпендикулярен касательной:

\[
OT\perp AB.
\]

Так как $\triangle AOB$ прямоугольный, отрезок $OT$ является высотой,
проведённой к гипотенузе $AB$.

Если

\[
AT=x,
\qquad
TB=y,
\]

то по теореме о высоте к гипотенузе

\[
OT^2
=
AT\cdot TB.
\]

Следовательно,

\[
r^2=xy,
\]

\[
r=\sqrt{xy}.
\]

Высота трапеции равна диаметру вписанной окружности:

\[
h=2r.
\]

Получаем

\formula{
h=2\sqrt{xy}
}

Например, если точка касания делит боковую сторону на отрезки $4$ и $1$,
то

\[
r=\sqrt{4\cdot1}=2,
\]

поэтому

\[
h=4.
\]

% =========================================================

\section*{6. Прямоугольный треугольник: три ключевые формулы}
\sectionline

Пусть в прямоугольном треугольнике $ABC$

\[
\angle ABC=90^\circ,
\]

а $BH$ --- высота к гипотенузе $AC$.

\begin{center}
\begin{tikzpicture}[scale=0.9]

\tkzSetUpLine[
  line width=1pt,
  color=accent
]

\tkzSetUpPoint[
  color=accent,
  fill=accent,
  size=3
]

\tkzSetUpLabel[
  color=black,
  font=\small
]

\tkzDefPoint(0,0){A}
\tkzDefPoint(7,0){C}
\tkzDefPoint(2,3.1623){B}
\tkzDefPoint(2,0){H}

\tkzDrawPolygon(A,B,C)
\tkzDrawSegment(B,H)

\tkzMarkRightAngle[size=0.28](A,B,C)
\tkzMarkRightAngle[size=0.28](B,H,C)

\tkzDrawPoints(A,B,C,H)

\tkzLabelPoints[below](A,H,C)
\tkzLabelPoints[above](B)

\end{tikzpicture}
\end{center}

Высота к гипотенузе:

\formula{
BH^2=AH\cdot HC
}

Первый катет:

\formula{
AB^2=AH\cdot AC
}

Второй катет:

\formula{
BC^2=HC\cdot AC
}

Полезно замечать структуру последних двух формул:

\[
\boxed{
\text{катет}^2
=
\text{гипотенуза}
\cdot
\text{проекция этого катета}
}
\]

Все три формулы следуют из подобия треугольников.

% =========================================================

\section*{7. Отрезок между серединами диагоналей}
\sectionline

Пусть основания трапеции $ABCD$ равны

\[
AD=b,
\qquad
BC=a,
\qquad
b>a.
\]

Пусть $M$ и $N$ --- середины диагоналей трапеции.

Отрезок $MN$ параллелен основаниям.

Его длина равна полуразности оснований:

\formula{
MN=\frac{b-a}{2}
}

Если заранее неизвестно, какое основание больше, используйте универсальную
запись

\formula{
MN=\frac{|b-a|}{2}
}

Почему появляется именно полуразность?

В соответствующих треугольниках возникают средние линии длиной

\[
\frac b2
\quad\text{и}\quad
\frac a2.
\]

Искомый отрезок равен разности этих длин:

\[
MN
=
\frac b2-\frac a2
=
\frac{b-a}{2}.
\]

% =========================================================

\section*{8. Как выбирать дополнительное построение}
\sectionline

В задачах о трапеции полезно придерживаться следующей стратегии.

\begin{enumerate}
  \item Сначала запишите формулу искомой величины.
  
  \item Определите, какого элемента не хватает.
  
  \item Если не хватает высоты, ищите треугольник, содержащий эту высоту.
  
  \item Если подходящего треугольника нет, рассмотрите два основных хода:
  \begin{itemize}
    \item проведение перпендикуляра;
    \item параллельный перенос боковой стороны или диагонали.
  \end{itemize}
  
  \item После построения выясните, что известно о полученном треугольнике:
  \begin{itemize}
    \item три стороны;
    \item прямой угол;
    \item равные стороны;
    \item угол $30^\circ$ или $60^\circ$;
    \item медиана;
    \item высота к гипотенузе.
  \end{itemize}
  
  \item Примените соответствующую теорию треугольников.
\end{enumerate}

\important{Главный ориентир.}
Задача формально может быть дана про трапецию, однако основная вычислительная
работа очень часто происходит внутри правильно выбранного треугольника.

% =========================================================

\section*{9. Типичные ошибки и самопроверка}
\sectionline

\begin{itemize}

\item
Не путайте полусумму оснований

\[
\frac{a+b}{2}
\]

и полуразность

\[
\frac{b-a}{2}.
\]

\item
После параллельного переноса обязательно объясняйте, почему соответствующие
отрезки равны.

\item
При применении формулы Герона сначала вычисляйте полупериметр.

\item
Проверяйте существование треугольника по неравенству треугольника.

\item
Равенство диагоналей является свойством равнобедренной трапеции.

\item
Перпендикулярность диагоналей должна быть дана или доказана отдельно.

\item
Если окружность вписана в трапецию, расстояние между основаниями равно
диаметру окружности:

\[
h=2r.
\]

\item
В задаче с точкой касания на боковой стороне сначала докажите или используйте
доказанный факт

\[
\angle AOB=90^\circ.
\]

\item
Для отрезка между серединами диагоналей используйте модуль разности оснований,
если неизвестно, какое основание больше:

\[
\frac{|a-b|}{2}.
\]

\item
Если в задаче требуется площадь, полезно в самом начале записать

\[
S=\frac{a+b}{2}h.
\]

Это сразу показывает, какого элемента не хватает.

\end{itemize}

% =========================================================
% Тренировочный блок
% =========================================================

\clearpage

\section*{Тренировочный блок. Путь А}
\sectionline

\textbf{Указание.}
Решайте задачи с полным обоснованием. Если выполняется дополнительное
построение, поясните его цель и укажите, какие равенства или свойства
из него следуют.

\begin{enumerate}

\item
Основания трапеции равны $10$ и $24$, боковые стороны равны $13$ и $15$.
Найдите площадь трапеции.

\item
В равнобедренной трапеции основания равны $8$ и $20$, боковая сторона
равна $10$. Найдите высоту и площадь трапеции.

\item
Диагональ равнобедренной трапеции равна $14$ и образует с основанием
угол $60^\circ$. Найдите среднюю линию трапеции.

\item
В трапецию вписана окружность. Точка касания окружности с одной из боковых
сторон делит эту сторону на отрезки длиной $4$ и $9$.
Найдите высоту трапеции.

\item
Основания трапеции равны $7$ и $19$.
Найдите длину отрезка, соединяющего середины её диагоналей.

\end{enumerate}

% =========================================================
% Ответы
% =========================================================

\clearpage

\section*{Ответы}
\sectionline

\renewcommand{\arraystretch}{1.35}

\begin{table}[ht]
\centering

\begin{tabularx}{\linewidth}{@{}cX@{}}

\toprule

\textbf{№}
&
\textbf{Ответ}
\\

\midrule

1
&
$204$
\\

2
&
$h=8$, \quad $S=112$
\\

3
&
$7$
\\

4
&
$12$
\\

5
&
$6$
\\

\bottomrule

\end{tabularx}

\end{table}

\end{document}